\section{Lefschetz for Picard groups in positive characteristic}\label{sec:A}

\subsection{The question and its quantifiers}
At 57:97 the transcription asks:
\begin{quote}
``Surjectivité (donc bijectivité) de $\Pic(X)\to\Pic(Y)$ pour $n\geq4$
sans conditions cohomologiques'';
\end{quote}
the adjacent box reads ``noyau nul si $n\geq4$''.
We take $X$ smooth, connected and projective over an algebraically closed
field $k$, $L$ a specified very ample invertible sheaf, and $Y=V(h)$ a
smooth member of $|L|$.  We distinguish an arbitrary smooth member from
the geometric generic member.  A tangent-space assertion for a Picard
scheme is not an assertion about the kernel on its $k$-points.

\paragraph{Current state of A1--A3.}
A3 is proved affirmatively by Theorem~\ref{a:very-ample-example}: an
explicit characteristic-two surface has a very ample sheaf with
one-dimensional negative first cohomology. The least requested dimension
is two. The unrestricted isomorphism A1 and injectivity A2 remain
undecided. The finite obstruction description A4 is proved.

\subsection{A finite, exact obstruction description}
Write $I=\OO_X(-Y)$ and $Y_m=V(I^{m+1})$, so $Y_0=Y$.
The cohomology of sheaves of units below is Zariski cohomology.

\begin{theorem}[A4: exact finite description; Proved]\label{a:finite}
Suppose $\dim X=n\geq4$.  Choose $N\geq0$ such that
\[
 H^i(Y,L^{-m}|_Y)=0\qquad(i=1,2;\ m>N).
\]
Such an $N$ exists.  On the topological space $Y$ put
\[
 U_N=\ker(\OO_{Y_N}^{\times}\longrightarrow\OO_Y^{\times})
     =1+I/I^{N+1}.
\]
For $N=0$ this is the zero group sheaf.  Restriction gives a canonical
identification $\Pic(X)\simeq\Pic(Y_N)$ and an exact sequence
\begin{equation}\label{a:exact}
 0\longrightarrow H^1(Y,U_N)\longrightarrow\Pic(X)
 \longrightarrow\Pic(Y)\xrightarrow{\delta_N}H^2(Y,U_N).
\end{equation}
Consequently the complete answers for this particular pair are
\[
 \ker(\Pic(X)\to\Pic(Y))=H^1(Y,U_N),\qquad
 \operatorname{coker}(\Pic(X)\to\Pic(Y))=\operatorname{im}(\delta_N).
\]
The sheaf $U_N$, with its actual multiplication, has the filtration
\[
 U_N^{(m)}=1+I^m/I^{N+1},\qquad
 U_N^{(m)}/U_N^{(m+1)}\simeq L^{-m}|_Y\quad(1\leq m\leq N).
\]
Thus a finite tower with the additive layers of R4 determines both groups
and all obstruction maps.  Replacing this tower by the direct sum of the
cohomology of its layers is not justified.
\end{theorem}
\begin{proof}
We first check algebraization.  The effective Lefschetz theorem
\cite[Expos\'e XII, Corollary 3.4, pp.~119--120]{SGA2} applies: closed
points of $Y$ have depth $n-1\geq2$, those of $X$ have depth $n\geq3$,
and $Y$ is an effective ample Cartier divisor.  Every line bundle on
$\widehat X_Y$ algebraizes on an open neighborhood $V$ of $Y$;
completion is fully faithful on vector bundles on $V$.
Every positive-dimensional projective closed subset meets $Y$: otherwise
$h$ would trivialize the ample sheaf $L$ on that subset.  Thus
$X\setminus V$ has codimension at least two.  Since $X$ is regular,
closing a divisor representing a line bundle on $V$ extends it uniquely
to $X$, by local factoriality.  We obtain
\begin{equation}\label{a:formal}
 \Pic(X)\simeq\Pic(\widehat X_Y).
\end{equation}

By \cite[Expos\'e XII, Corollary 3.5, p.~120]{SGA2}, $Y$ is connected.
Hence $H^0(Y,\OO_Y)=k$.  The anti-ample sheaves $L^{-m}|_Y$ have no
nonzero section for $m>0$.  Induction using
$0\to L^{-m}|_Y\to\OO_{Y_m}\to\OO_{Y_{m-1}}\to0$ gives
$H^0(Y_m,\OO_{Y_m})=k$: its injection into $k$ contains the constants.
All transition maps on global units are consequently the identity on
$k^\times$.

A formal line bundle is a compatible system of line bundles on the $Y_m$.
Compatible isomorphism classes can be equipped successively with
compatible isomorphisms.  The ambiguity is the derived inverse limit of
their global automorphism groups; it vanishes because the transition maps
on $k^\times$ are surjective.  Therefore
$\Pic(\widehat X_Y)=\varprojlim_m\Pic(Y_m)$.
For $m\geq1$, multiplication on $1+I^m/I^{m+1}$ is addition, and the
calculation of global units gives
\begin{equation}\label{a:layer}
 0\to H^1(Y,L^{-m}|_Y)\to\Pic(Y_m)\to\Pic(Y_{m-1})
       \xrightarrow{o_m}H^2(Y,L^{-m}|_Y).
\end{equation}
For $m>N$ its middle transition is an isomorphism, so the inverse limit
is $\Pic(Y_N)$.  The bound exists by Serre duality and Serre vanishing
on the $(n-1)$-fold $Y$, since $(n-1)-i>0$ for $i=1,2$.
Finally the cohomology sequence of
$1\to U_N\to\OO_{Y_N}^{\times}\to\OO_Y^{\times}\to1$ and
surjectivity on global units give \eqref{a:exact}.  Multiplication of units
gives the stated filtration.  These are the actual restriction maps of
the given thickening, so no extension or obstruction has been discarded.
\end{proof}

\paragraph{An explicit finite obstruction procedure.}
Choose a finite affine cover of $Y_N$ trivializing a line bundle $M$ on
$Y$.  Lift its transition units $g_{ij}$ to units $\widetilde g_{ij}$ on
the pairwise intersections in $Y_N$.  The products
\[
 c_{ijk}=\widetilde g_{ij}\widetilde g_{jk}\widetilde g_{ki}\in U_N
\]
represent $\delta_N(M)$.  Different lifts change it by a coboundary;
$M$ extends if and only if this class vanishes.  Alternatively solve the
equations first modulo $I^2$, then modulo $I^3$, through $I^{N+1}$.
At stage $m$ the remaining additive cocycle lies in $L^{-m}|_Y$ and
represents $o_m$; all choices of a lift form a torsor under
$H^1(Y,L^{-m}|_Y)$.  Those choices must be retained because the next
obstruction can depend on them.  Quasi-coherent acyclicity on affine
intersections, applied to the finite filtration, justifies these
\v Cech computations.  This is a finite description using the layers and
their multiplication; the vector spaces alone do not include its maps.

\begin{corollary}[A4b: bounded finite $p$-groups; Proved]\label{a:pgroups}
In characteristic $p>0$, for $n\geq4$, both groups in
Statement~\ref{a:finite} are finite and killed by $p^e$ whenever
$p^e\geq N+1$.  For $n=3$ the kernel is still finite and killed by a
power of $p$.
\end{corollary}
\begin{proof}
For $a\in I/I^{N+1}$, $(1+a)^{p^e}=1+a^{p^e}=1$.  Thus $p^e$ kills
$U_N$, its cohomology, and the two groups in \eqref{a:exact}.
The $p^e$-torsion of $\Pic(X)$ is finite: $\NS(X)$ is finitely generated,
and the $k$-points of $\mathbf{Pic}^0_{X/k}$ are those of an abelian
variety, with finite multiplication-by-$p^e$ kernel.  This proves
finiteness of the kernel.  Multiplication by $p^e$ is surjective on the
$k$-points of an abelian variety.  Since the cokernel is killed by $p^e$,
the image contains $\Pic^0(Y)(k)$.  It is therefore a bounded torsion
quotient of the finitely generated $\NS(Y)$ and is finite.

For $n=3$, choose $N$ so that only $H^1(Y,L^{-m}|_Y)=0$ for $m>N$.
The transitions in \eqref{a:layer} are then injective.
The identification \eqref{a:formal} still holds and embeds $\Pic(X)$
into $\Pic(Y_N)$.  Its restriction kernel embeds in $H^1(Y,U_N)$;
the exponent and finiteness arguments apply unchanged.  No surjectivity
of the transitions has been asserted in dimension three.
\end{proof}

\begin{corollary}[A4c: vanishing criterion; Known, hypotheses checked]
\label{a:vanishing}
If $n\geq3$ and $H^1(Y,L^{-m}|_Y)=0$ for all $m>0$, restriction
$\Pic(X)\to\Pic(Y)$ is injective.  If $n\geq4$ and also all the
groups $H^2(Y,L^{-m}|_Y)$ vanish, restriction is an isomorphism.
\end{corollary}
\begin{proof}
Each transition in \eqref{a:layer} is injective in the first case and an
isomorphism in the second; we use \eqref{a:formal}.  This also checks
the smooth case of \cite[Expos\'e XII, Corollary 3.6, p.~121]{SGA2}:
its local depth and parafactoriality conditions follow from regularity.
The vanishing assumptions are absent from the unrestricted question A1.
\end{proof}

\subsection{Reduction of the main gates}
\begin{proposition}[A3r: reduction to surfaces; Proved]\label{a:surface}
If there is a smooth projective variety of dimension at least two with
a very ample $L$ and $H^1(X,L^{-1})\ne0$, there is such a surface.
If A3 has a positive existence answer, its least dimension is two.
\end{proposition}
\begin{proof}
When $\dim X>2$, Serre duality and Serre vanishing give
$H^1(X,L^{-a})=0$ for large $a$.  Choose the largest positive $a$ for
which it is nonzero, and a smooth integral $D\in|L|$ by Bertini for
the specified embedding.  The sequence
\[
 0\to L^{-a-1}\to L^{-a}\to L^{-a}|_D\to0
\]
gives $0\ne H^1(X,L^{-a})\hookrightarrow H^1(D,L^{-a}|_D)$.
The sheaf $L^a|_D$ is very ample.  Repeat until the dimension is two.
This proves an implication between existence questions; it does not
provide the initial example.
\end{proof}

\begin{proposition}[A1r: reduction to fourfolds; Proved]\label{a:fourfold}
The universal assertion A1 is equivalent to its dimension-four case.
A failure of injectivity in any dimension $n\geq4$ also gives one in
dimension four.
\end{proposition}
\begin{proof}
For $n>4$, choose large $d$ and a general smooth $D\in|L^d|$ transverse
to $Y$.  Serre duality and vanishing on $X$ give
\[
 H^j(X,L^{-a})=0\qquad(j=1,2,3;\ a\geq d)
\]
after increasing $d$.  The restriction sequences with exponents $md$
and $(m+1)d$ give, for every $m\geq1$,
\[
 H^1(D,L^{-md}|_D)=H^2(D,L^{-md}|_D)=0.
\]  The same argument on $Y$, whose dimension is at least four,
gives both vanishings on $Y\cap D$ for its normal powers in $Y$.
Corollary~\ref{a:vanishing} gives vertical isomorphisms in
\[
 \begin{matrix}
 \Pic(X)&\longrightarrow&\Pic(Y)\\
 \big\downarrow\wr&&\big\downarrow\wr\\
 \Pic(D)&\longrightarrow&\Pic(Y\cap D).
 \end{matrix}
\]
Here $Y\cap D$ is a smooth hyperplane for $L|_D$, the restriction of the
\emph{original} embedding.  Only the auxiliary cut used $L^d$.
Both rows have the same kernel and cokernel.  Iteration proves the claim.
\end{proof}

\begin{proposition}[A2b: dimension two fails; Proved]\label{a:segre}
In every characteristic, a smooth hyperplane section of a smooth
projective surface can have a nonzero restriction kernel.
\end{proposition}
\begin{proof}
Let $X=\PP^1\times\PP^1$ with the Segre embedding and
$L=\OO(1,1)$.  The diagonal is the smooth hyperplane section
$x_0y_1-x_1y_0=0$, including in characteristic two.
The map $\Pic(X)=\ZZ^2\to\Pic(Y)=\ZZ$ is $(a,b)\mapsto a+b$;
its kernel is $\ZZ(1,-1)$.  A universal injectivity threshold, if it
exists, is at least three.  No dimension-three or dimension-four
injectivity assertion follows from this example.
\end{proof}

\subsection{The infinitesimal kernel on 57:100}
The supplied transcription calls this kernel a ``$p$-groupe fini
unipotent, with the marginal qualification ``purement infinitésimal
(57:100). Its phrase ``sans partie fixe must also be retained: the
map on classes defined over $k$ is distinguished from the group-scheme
map after geometric-generic base change.
\begin{proposition}[R3/A3t: tangent calculation; Proved]\label{a:tangent}
For $\dim X\geq2$ and any smooth hyperplane $Y$, the scheme-theoretic
restriction kernel $K$ satisfies
\[
 \Lie K=\ker\bigl(H^1(X,\OO_X)\to H^1(Y,\OO_Y)\bigr)
        =h\,H^1(X,L^{-1}).
\]
Multiplication by $h$ on this last group is injective; therefore
$\dim_k\Lie K=h^1(X,L^{-1})$.  If the kernel under consideration is
finite and connected, it is trivial exactly when this group vanishes.
\end{proposition}
\begin{proof}
First-order line-bundle deformations have transition units
$1+\epsilon a_{ij}$, giving tangent space $H^1(\OO)$; tangent spaces
preserve kernels.  Apply the long exact sequence of
$0\to L^{-1}\xrightarrow{h}\OO_X\to\OO_Y\to0$, observing that
$H^0(X,\OO_X)\to H^0(Y,\OO_Y)$ is the identity of $k$.
For a finite connected group scheme over $k$, vanishing of the tangent
space gives $\mathfrak m/\mathfrak m^2=0$ in its local Artinian ring.
Nakayama's lemma gives $\mathfrak m=0$, so it is trivial.
\end{proof}

\begin{proposition}[R3g: geometric generic kernel; Proved]\label{a:generic}
For the geometric generic hyperplane, the kernel on $\mathbf{Pic}^\tau$
is finite and connected.  Thus the equivalence in R3 is valid for smooth
$X$: this kernel is nontrivial if and only if $H^1(X,L^{-1})\ne0$.
\end{proposition}
\begin{proof}
We first check finiteness without presupposing the manuscript's assertion.
For normal projective $X$, $\mathbf{Pic}^\tau$ is proper and of finite
type: its identity component is proper
\cite[Theorem 5.4 and Remark 5.6]{Kleiman2005}, and its component group
is the finite torsion subgroup of $\NS(X)$.  The restriction kernel is
therefore a proper group scheme of finite type.  If its reduced identity
component had positive dimension, it would be a positive-dimensional
abelian variety and would have nonzero $\ell$-torsion points for a prime
$\ell\ne\operatorname{char}k$.  Such a point is impossible: Kummer
theory identifies prime-to-characteristic torsion line bundles with
$H^1(\mu_\ell)$, and the surjectivity of the fundamental-group map for
an ample divisor gives injectivity of restriction on that group
\cite[Expos\'e XII, Corollary 3.5]{SGA2}.  Thus the kernel is finite.

We prove that it has no nonzero geometric point.  Every point of a finite
group has finite order.  For each integer $a>0$, the $a$-torsion points
of $\Pic^\tau(X)$ are finite and defined over the algebraically closed
field $k$.  Hence every point in the geometric generic kernel comes from
a torsion line bundle on $X$.  Prime-to-$p$ torsion has already been
excluded.  To exclude $p$-torsion, we use the injective map
\[
 d\log:\Pic(X)[p]\hookrightarrow H^0(X,\Omega_X^1).
\]
Indeed, trivialize the $p$th power of a line bundle, obtaining local
units $u_i$ whose ratios are $p$th powers.  The forms $du_i/u_i$ glue.
If their common form is zero, then each $u_i$ has a $p$th root in the
function field, since $k$ is perfect.  Normality makes these roots local
units, and they trivialize the bundle.  The construction commutes with
restriction to a smooth divisor.

No nonzero form defined over $k$ restricts to zero on the geometric generic
hyperplane.  To see this, use the incidence variety of pairs $(x,H)$ with
$x\in X\cap H$.  On the locus where the hyperplane is smooth at $x$,
the relative cotangent sheaf is locally free.  Vanishing on its generic
fibre would make the pulled-back form zero throughout this integral
smooth locus.  But, at a point $x$ where the form is nonzero, the tangent
spaces of hyperplanes through $x$ run through all codimension-one
subspaces of $T_xX$.  Since $\dim X\geq2$, some such subspace is not
annihilated by the form.  This is a contradiction.  Thus $p$-torsion is
excluded; successively taking $p$th powers excludes $p^a$-torsion for
all $a$.  The finite kernel has just its identity as a geometric point
and is connected.  Proposition~\ref{a:tangent} finishes the proof.
\end{proof}

This argument checks the finite-infinitesimal property actually needed
for R3; it does not need the stronger unipotence assertion written at
57:100.  Its incidence argument is only injectivity for forms defined
over $k$.  We used it for torsion points that have first been proved to
be defined over $k$; we did not incorrectly infer injectivity for all
linear combinations with geometric generic coefficients.

\paragraph{A transcription correction, kept explicit.}
The two displayed identifications on 57:100 both print
$\mathrm H^1(X,\underline{\mathcal O}_X)^{1-F}$, including in
the $\mu_p$ line. In the additive line the correct expression is
$\ker(F:H^1(\OO_X)\to H^1(\OO_X))$ for proper connected
$X$ over the perfect field; $1-F$ instead belongs to the
Artin--Schreier sequence for $\mathbf Z/p$. The second line also
writes the quotient $\mathcal O_X/\mathcal O_X^*$, whereas the
multiplicative quotient needed there is
$\mathcal O_X^\times/\mathcal O_X^{\times p}$.  As printed in the supplied
transcription it is not correct.  For smooth connected proper $X/k$ the
Kummer sequence gives $H^1_{\mathrm{fppf}}(X,\mu_p)=\Pic(X)[p]$,
since $k^\times$ is $p$-divisible.  The corresponding Zariski quotient
is $\OO_X^\times/\OO_X^{\times p}$, and its map is $d\log$.
We use this corrected mathematical statement, not a silent change of
the diplomatic text.

\subsection{A sharper surface gate for A3}
\begin{proposition}[A3n: a necessary Frobenius-nilpotence bound; Known deduction]
\label{a:nilpotence}
For a smooth projective connected variety $X/k$ of dimension at least two,
let $q_{\mathrm{nil}}(X)$ be the dimension of the subspace of
$H^1(X,\OO_X)$ killed by some iterate of absolute Frobenius.
If $L$ is very ample and $H^1(X,L^{-1})\ne0$, then
\[
 q_{\mathrm{nil}}(X)\geq h^0(X,L).
\]
In the surface case we furthermore have $h^0(X,L)\geq5$.
Consequently A3 reduces to smooth surfaces with
$q_{\mathrm{nil}}(X)\geq h^0(X,L)\geq5$.
\end{proposition}
\begin{proof}
Apply \cite[Theorem 4.20, pp.~15--16]{Enokizono2021} with the
integral ample divisor $D$ of $L$, $M=D$, and $E=0$.
The variety is normal and geometrically connected over an infinite
perfect field. General hyperplanes cut a smooth surface by Bertini;
$M$ remains ample, hence nef and big, on that surface. The theorem
therefore gives $H^1(X,L^{-1})=0$ whenever
$h^0(X,L)-1\geq q_{\mathrm{nil}}(X)$. Its contrapositive is the
first assertion. In dimension two, a complete very ample system
with $h^0\leq4$ embeds the surface as $\PP^2$ or a hypersurface
in $\PP^3$. The usual hypersurface exact sequence and the
cohomology of line bundles on projective space give
$H^1(X,\OO_X)=0$ in both cases, contrary to the first assertion.
Combining with Proposition~\ref{a:surface} proves the last claim.
\end{proof}
The theorem applies in characteristics two and three as well.
It rules out ordinary varieties, in the explicit sense that Frobenius
is injective on $H^1(\OO_X)$, but supplies no existence theorem for
the remaining surfaces. Its Frobenius-nilpotent dimension must not
be confused with the nonreduced dimension of the Picard scheme.

\begin{proposition}[A3s: general type and numerical bounds; Known deduction]
\label{a:generaltype}
Let $X/k$ be a smooth projective surface and let $L$ be very ample.
If $H^1(X,L^{-1})\ne0$, then $X$ is of general type. For its
smooth minimal model $X_0$, we have
\[
 L^2\leq
 \begin{cases}
 K_{X_0}^2,&p\geq3,\\
 \max\{K_{X_0}^2,K_{X_0}^2-3\chi(\OO_X)+2\},&p=2.
 \end{cases}
\]
Moreover $X$ does not admit a smooth lifting over $W_2(k)$.
\end{proposition}
\begin{proof}
If $X$ is neither of general type nor quasi-elliptic of Kodaira
dimension one, \cite[Theorem 1.6(1), pp.~191--192]{Terakawa1999}
gives the required vanishing for every nef and big $L$, hence for
our very ample $L$. In the remaining quasi-elliptic case let $F$
be a general fibre. This is an integral Gorenstein curve of arithmetic
genus one with $\omega_F\simeq\OO_F$. A positive-degree invertible
sheaf $M$ on $F$ has $h^0(F,M)=\deg M$ by duality and Riemann--Roch.
In degrees one and two its complete system cannot embed $F$:
it maps to a point or to $\PP^1$. Thus very ampleness of $L|_F$
forces $L\cdot F\geq3$. The corrected vanishing theorem
\cite[Theorem 3.1, pp.~4--5]{Zhang2022} applies: its bounds are
$L\cdot F>2$ for $p=2$ and $L\cdot F>1$ for $p=3$.
This excludes the quasi-elliptic case. The two numerical bounds
are contrapositives of \cite[Theorem 1.6(2)--(3)]{Terakawa1999}.
Finally, in the presence of a smooth $W_2(k)$ lifting,
\cite[section 11, Corollary 11.2]{Langer2011} with the numerically
flat bundle $E=\OO_X$, $i=0$, and $j=1$ gives vanishing because
$1<\min(p,2)$ for every prime $p$. This proves the last assertion.
\end{proof}
We use Zhang's corrected bounds rather than the stronger earlier
quasi-elliptic assertions which that paper disproves in characteristic two.
Combining Propositions~\ref{a:surface}, \ref{a:nilpotence}, and
\ref{a:generaltype}, any A3 example must be a surface
of general type with a very ample $L$, the displayed numerical bounds,
$q_{\mathrm{nil}}\geq h^0(L)\geq5$, and no smooth $W_2$ lifting.
These were necessary restrictions during the search; the explicit
construction below now proves existence.

\subsection{Tested approaches and their exact barriers}
\begin{proposition}[A3f: the indicated Raynaud polarizations are not very ample;
Proved]\label{a:fibre}
Let $p$ be the characteristic and let $\ell\geq2$ be prime to $p$.
Let $F$ be the projective curve obtained from
$\Spec k[u^p,u^\ell]$ by adding its smooth point $\infty$, with
normalization $\PP^1_u$.  For $0\leq a<\ell$, the complete system
$|\OO_F(a\infty)|$ does not give an immersion on a dense open subset
of $F$.  Consequently a line bundle restricting to this sheaf on such
a fibre cannot be very ample on the ambient variety.
\end{proposition}
\begin{proof}
A section is a polynomial in the monomials of the semigroup
$\langle p,\ell\rangle$ whose pole order at infinity is at most $a$.
Because $a<\ell$, every such exponent is a multiple of $p$.
Thus every section is a polynomial in $u^p$.  On the smooth locus
$u\ne0$, any map defined by these sections has zero differential, or
is constant.  It cannot be an immersion.  Restriction of a closed
immersion to a closed fibre would still be a closed immersion, proving
the last assertion.
\end{proof}
This applies to the polarizations $\OO(a\widetilde E)\otimes\phi^*N^b$
in the nonvanishing range $a\leq\ell-1$ of
\cite[Theorem 3.7]{Zheng2015}: on a geometric fibre the pullback factor
is trivial and $\widetilde E$ is its smooth point at infinity.  It also
applies to $\widetilde D=(\ell-1)F_\infty+g^*D'$ in
\cite[Definition 2.1 and Proposition 2.3, pp.~7--8]{Mukai2013revised},
where the fibre is explicitly $T^\ell=S^p$ and is normalized at infinity.
We have therefore closed the direct very-ampleness route for these
particular stated polarizations, not merely failed to find an embedding.
The finite semigroup checks in \texttt{code/check\_raynaud\_fibre.py}
cover 518 pairs with primes at most 19 and $2\leq\ell\leq80$;
the preceding proof covers all pairs.

The finite $p$-group assertion of R4 is proved above; its vanishing is
not.  We have checked the following specific limitations.
\begin{enumerate}[leftmargin=*]
\item The classical Picard theorem needs the vanishings in
Corollary~\ref{a:vanishing}; the editorial note in
\cite[Expos\'e XII, section 4, note 5, p.~128]{SGA2} explicitly separates
the characteristic-$p$ issue.
\item \cite[Theorem 0.1]{BhattDeJong2014} gives local injectivity up
to $p$-power torsion, not an isomorphism up to that torsion.  Thus the
recollection in SA14-H3 requires correction.  Its local examples do not
give the smooth projective pair required here.
\item The results of Cotner--Zavyalov
\cite[Theorem 3.4.4 and Corollary 3.4.5]{CotnerZavyalov2024}
require strong ampleness or additional cohomological hypotheses.
Examples 3.3.6 and 3.4.6 distinguish failure for the Picard scheme
from failure for its group of line bundles.  Tangent noninjectivity alone
does not settle A1 or A2.
\item The very ample family in \cite[section 2]{LauritzenRao1997}
has first nonzero negative cohomology in degree $3r-4$ on a
$(3r-3)$-fold, $r\geq3$; all lower degrees vanish by the same
calculation.  It does not supply the degree-one nonvanishing in A3.
\item Proposition~\ref{a:fibre} shows that the indicated Raynaud-type
nonvanishing polarizations are not very ample.  Passing to an unspecified
large power can remove the nonvanishing and is not a valid solution.
\end{enumerate}
For A1 the current gate is a smooth fourfold with a nonzero finite
obstruction in Statement~\ref{a:finite}, or a proof that every such
obstruction vanishes. For A2 it is the surviving $p$-torsion kernel and
the universal dimension-three case. These two gates remain unfinished.
The following construction settles the separate existence question A3.

\subsection{A3R: exclusion of every positive two-parameter Raynaud polarization}
\begin{lemma}[Connection bound; Proved]\label{a:connection-bound}
Let $C$ be a smooth projective connected curve over an algebraically closed
field of characteristic $p>0$. Let $L$ have positive degree, and let
$0\to\OO_C\to E\to L\to0$ be nonsplit with split Frobenius pullback.
For $0\leq n<p$ and any line bundle $M$ with
$\deg M<n\deg L$, we have $\Hom_C(\operatorname{Sym}^n E,M)=0$.
\end{lemma}
\begin{proof}
The splitting $F^*E=\OO_C\oplus L^p$ is unique since
$\Hom(L^p,\OO_C)=0$. Its canonical connection preserves $\OO_C$.
The second fundamental map $\beta:L^p\to\Omega_C^1$ is nonzero:
otherwise the splitting is horizontal and Cartier descent splits the
original extension. The precise descent source is Katz, \emph{Nilpotent
connections and the monodromy theorem}, Theorem~5.1, printed pp.~190--191 (opened in this run).
For a homomorphism $\varphi:\operatorname{Sym}^nE\to M$, its Frobenius
pullback is horizontal for the canonical connections. Its component on
$L^{pn}$ is zero by the degree inequality. On the symmetric power, the
connection carries the summand of exponent $j$ to the summand of exponent
$j-1$ with coefficient $j\beta$, in addition to terms of the same exponent.
If its components of exponents at least $j$ are zero, horizontality forces
$j\beta$ times its component of exponent $j-1$ to vanish. Since
$1\leq j\leq n<p$ and $\beta$ is nonzero, that component vanishes too.
Induction gives $F^*\varphi=0$, hence $\varphi=0$ by faithful flatness of
Frobenius on the smooth curve. For $n=0$ the assertion is simply that a
negative-degree line bundle has no sections.
\end{proof}

\begin{proposition}[A3R; Proved]\label{a:raynaud-all}
For the smooth Raynaud cyclic-cover surface of Zheng's construction,
with $\ell\mid p+1$, $L=N^\ell$ on its Tango base curve and
$\deg N>0$, let $\widetilde E$ be the positive ramification section.
For every $a,b>0$, the line bundle
$Z_{a,b}=\OO_X(a\widetilde E)\otimes\phi^*N^b$ cannot be both very
ample and have $H^1(X,Z_{a,b}^{-1})\ne0$.
\end{proposition}
\begin{proof}
Set $d=(p+1)/\ell$. For $a=\ell$, the finite-cover pushforward in
Zheng's Proposition~3.3 is
$\bigoplus_{i=0}^{\ell-1}\OO_{\PP(E)}(-id-1)\otimes\pi^*N^{ip-b}$.
Its direct image to $C$ is zero. Relative duality for the quotient
convention $\pi_*\OO_{\PP(E)}(1)=E$ therefore gives
\[
 H^1(X,Z_{\ell,b}^{-1})=
 \bigoplus_{i=1}^{\ell-1}
 H^0\!\left(C,(\operatorname{Sym}^{id-1}E)^\vee
                         \otimes N^{ip-b-\ell}\right).
\]
Here $n=id-1$ lies between zero and $p-d$, hence below $p$, and
\[
 \deg(N^{ip-b-\ell})-n\deg L=-(i+b)\deg N<0.
\]
The lemma makes every summand zero.
Next $\widetilde E^2=\deg N>0$, and $\phi|_{\widetilde E}$ is the
identity of $C$. Thus $Z_{a,b}^{-1}|_{\widetilde E}$ has negative
degree $-(a+b)\deg N$. The divisor exact sequence consequently
injects $H^1(X,Z_{a+1,b}^{-1})$ into $H^1(X,Z_{a,b}^{-1})$.
We obtain vanishing for every $a\geq\ell$ and $b>0$.
For $a<\ell$, restriction to the cuspidal fibre is
$\OO_F(a\infty)$ and Proposition~\ref{a:fibre} rules out very
ampleness. These two cases exhaust all positive $a,b$.
\end{proof}

This excludes the entire displayed two-parameter family. It is not a
nonexistence theorem for other line bundles or other surfaces.
The connection argument and all inequalities received an independent
targeted check in this run. The exact Cech computation covers1008 cases;
it serves as a check on this proof, not as its substitute.

\subsection{A3: an explicit very ample surface with negative first cohomology}

\begin{theorem}[A3; Proved]\label{a:very-ample-example}
There is a smooth geometrically integral projective surface $Y$ over an
algebraically closed field of characteristic two, and a very ample
invertible sheaf $L$ on $Y$, for which
\[
 h^0(Y,L)=562,\qquad h^1(Y,L^{-1})=1,\qquad h^1(Y,\OO_Y)=565.
\]
Thus the least dimension in the range requested in A3, namely
$\dim X\geq2$, is two. For the geometric generic hyperplane of this
embedding the kernel on $\mathbf{Pic}^{\tau}$ is nontrivial and has
one-dimensional tangent space. This conclusion concerns the group
scheme kernel; it does not give a counterexample on geometric points
to A1 or A2.
\end{theorem}

We give all the defining data, including a finite algorithm specifying
every coefficient of the sections used in the construction. The final
ground field is an explicitly specified algebraic closure of a purely
transcendental extension of $\mathbb F_2$. No numerical search for smooth
branch divisors is substituted for the proof of smoothness.

\paragraph{1. The curve and the extension.}
Let $C$ be the smooth projective model over $\mathbb F_2$ of
\[
 y^2+y=x^{15}.
\]
The Artin--Schreier polynomial is irreducible over
$\overline{\mathbb F}_2(x)$: a rational function whose Artin--Schreier
difference has a pole has an even pole order there, whereas $x^{15}$
has an odd pole order at infinity. There is one point $\infty$ over
infinity, with $v_\infty(x)=-2$ and $v_\infty(y)=-15$.
Write $P_0=(0,0)$, $P_1=(0,1)$, and put $t=y/x^8$. Then
\[
 \operatorname{div}(t)=7P_0+\infty-8P_1,\qquad
 x+t^{-2}=t^{-1}x^{-7},\qquad dt=x^6\,dx.
\]
In particular $t$ is a uniformizer at $\infty$, and
$v_\infty(x+t^{-2})=13$. The affine curve is smooth and $dx$ is
nowhere zero there, since the derivative of its equation with respect
to $y$ is one. The displayed formula for $dt$ therefore gives
\[
 \operatorname{div}(dx)=12\infty,\qquad g(C)=7.
\]
These formulas also verify the extension and smoothness calculations
below directly, without an appeal to an unspecified Tango curve.

Set $N=\OO_C(2\infty)$ and $D=6\infty$, so that
$\OO_C(D)=N^3$. Cover $C$ by
$U=C\setminus\{\infty\}$ and $V=C\setminus\{P_0,P_1\}$.
Define a rank-two bundle $E$ by frames $(e_0,e_1)$ on $U$ and
$(v_0,v_1)$ on $V$, with transition
\begin{equation}\label{a:explicit-transition}
 v_0=e_0,\qquad v_1=t^{-6}(e_1-t^{-1}e_0).
\end{equation}
It fits into $0\to\OO_C\to E\to N^3\to0$.
The extension is nonsplit. Indeed a splitting would express $t^{-1}$
as a function regular on $U$ plus a function divisible by $t^6$ and
regular at $\infty$. The function on $U$ would have its sole pole,
of order one, at $\infty$. This is impossible: the ring of $U$ is
$\mathbb F_2[x,y]/(y^2+y-x^{15})$, and its possible pole orders are
$2a$ and $15+2a$, $a\geq0$.

The Frobenius pullback splits. On $U$ its second summand is generated
by $e_1^2+xe_0^2$; on $V$ the corresponding expression is
\[
 v_1^2+t^{-12}(x+t^{-2})v_0^2.
\]
The coefficient in this expression is regular on $V$ and has order
one at $\infty$. This gives a splitting of the quotient $N^6$ and
also a map $N^6\to\operatorname{Sym}^2E$.

\paragraph{2. The smooth cyclic-cover surface.}
Use the quotient convention $P=\mathbf P_C(E)$, so that
$\pi_*\OO_P(j)=\operatorname{Sym}^jE$ for $j\geq0$.
Let $E_\infty$ be the section $e_0=0$; its divisor sheaf is
$\OO_P(1)$ and $E_\infty^2=6$. The map just constructed defines
a divisor
\[
 C''\in\big|\OO_P(2)\otimes\pi^*N^{-6}\big|.
\]
On $U$, in the affine chart $e_0\ne0$ with $w=e_1/e_0$, its
equation is $w^2+x=0$. It is disjoint from $E_\infty$.
It is smooth on this chart since $dx$ never vanishes on $U$.
Over $\infty$ its equation is
$w_V^2+t^{-12}(x+t^{-2})=0$; the second term has order one,
so the same smoothness assertion holds there. Thus $C''$ and
$E_\infty$ are disjoint smooth divisors.

Put $M=\OO_P(-1)\otimes\pi^*N^2$. Then
$M^{-3}=\OO_P(E_\infty+C'')$. Let $\sigma$ be the product of
the two defining sections, and form the cyclic cover
\[
 \psi:S=\operatorname{Spec}_P
       (\OO_P\oplus M\oplus M^2)\longrightarrow P,
 \qquad M^3\xrightarrow{\sigma}\OO_P.
\]
This is finite flat of degree three. Locally at the branch locus it
has equation $z^3=u$ with $u$ a regular parameter. Since three is
invertible in characteristic two and the branch components are smooth
and disjoint, $S$ is smooth. It is projective because it is finite over
$P$. Moreover $H^0(S,\OO_S)=\mathbb F_2$: the two nonconstant
summands have negative degree on the fibres of $\pi$. This remains
true after extension of the ground field, and smoothness implies
geometric integrality. Write $f=\pi\psi$ and let $T$ be the
ramification section over $E_\infty$. Then
\[
 \psi^*E_\infty=3T,\qquad T\simeq C,\qquad T^2=2.
\]
On the dense affine chart used above all equations are simply
\begin{equation}\label{a:surface-affine}
 y^2+y=x^{15},\qquad z^3=w^2+x.
\end{equation}

\paragraph{3. An ample generated bundle with the required cohomology.}
Define
\[
 Q=\OO_S(2T)\otimes f^*N.
\]
The bundle $E$ is nef: after pullback to any smooth projective curve,
a quotient line bundle receives either a nonzero map from $\OO$ or
factors through the positive-degree quotient $N^3$; in either case
its degree is nonnegative. Thus $\OO_P(1)$ is nef. The line bundle
$\OO_P(2)\otimes\pi^*N^3$ has positive intersection with every
curve and positive self-intersection: on a fibre its degree is two,
and on a horizontal curve its base factor has positive degree.
The Nakai criterion makes it ample. Since
\[
 Q^3=\psi^*(\OO_P(2)\otimes\pi^*N^3),
\]
$Q$ is ample as well.

For an integer $a$, the local equation $z^3=u$ at $T$ gives
\begin{equation}\label{a:cyclic-floor}
 \psi_*\OO_S(aT)=
 \bigoplus_{i=0}^2 M^i\bigl(\lfloor(a+i)/3\rfloor E_\infty\bigr).
\end{equation}
Indeed the $i$th eigensummand consists locally of
$z^i u^{-\lfloor(a+i)/3\rfloor}$ times regular functions.
For $a=2$, this formula and projection to $C$ give
\[
 f_*Q=N\oplus N^3.
\]
Both summands are generated by global sections: their bases are
$(1,x)$ and $(1,x,x^2,x^3)$, respectively. The first summand generates
$Q$ away from $T$. The second has valuation $1-3+2=0$ along $T$
and generates there. Hence $Q$ is globally generated and
\[
 h^0(S,Q)=6;\qquad
 1,x,z,zx,zx^2,zx^3
\]
are its sections in the affine trivialization of
\eqref{a:surface-affine}.

For $a=-2$, formula~\eqref{a:cyclic-floor} instead gives
\[
 \psi_*Q^{-1}=
 \bigl(\OO_P(-1)\otimes\pi^*N^{-1}\bigr)
 \oplus\bigl(\OO_P(-2)\otimes\pi^*N\bigr)
 \oplus\bigl(\OO_P(-2)\otimes\pi^*N^3\bigr).
\]
Relative duality, with the convention stated above, says
$R^1\pi_*\OO_P(-2)=(\det E)^{-1}=N^{-3}$.
All zeroth direct images in the preceding display vanish, so Leray
gives
\begin{equation}\label{a:Q-cohomology}
 H^1(S,Q^{-1})=H^0(C,N^{-2}\oplus\OO_C)=\mathbb F_2.
\end{equation}
Proposition~\ref{a:raynaud-all}, applied with $a=2n$, $b=n$,
$p=2$ and $\ell=3$, also gives
\[
 H^1(S,Q^{-n})=0\quad(n\geq2).
\]
For later dimension checks, $R^1f_*\OO_S=N$, whence
$h^1(S,\OO_S)=7+h^0(C,N)=9$. Intersection and duality give
$Q^2=16$, $K_S=f^*N^5$, $K_SQ=20$, and $\chi(\OO_S)=-2$.
The canonical formula follows as well from
$K_P=\OO_P(-2)\otimes\pi^*(K_C\otimes\det E)$ and
$K_S=\psi^*(K_P\otimes M^{-2})$.

\paragraph{4. A specified very ample power and its complete basis.}
Set $B=\OO_P(6)\otimes\pi^*N^9$. Then
\[
 Q^9=\psi^*B,\qquad
 \psi_*Q^9=B\oplus BM\oplus BM^2
   =\bigoplus_{i=0}^2\OO_P(6-i)\otimes\pi^*N^{9+2i}.
\]
The bundle $\operatorname{Sym}^a E\otimes N^b$ has a filtration
with quotients $N^{b+3j}$, $0\leq j\leq a$.
For $(a,b)=(6,9),(5,11),(4,13)$ all these quotients, even after
subtracting any effective divisor of degree two on $C$, have degree
at least sixteen, greater than $2g(C)-2=12$.
Thus their first cohomology vanishes, as does that of the filtered
bundle. Restriction of sections is consequently surjective over every
base subscheme of length at most two.

It follows that $B$ is very ample on $P$. To see every infinitesimal
case explicitly, a length-two subscheme of $P$ either lies over a
length-two subscheme of $C$, or lies in one fibre. In the first case
use the just-proved surjectivity and relative very ampleness of
$\OO_P(6)$; in the second, use $\OO_{\mathbf P^1}(6)$.
The same argument with length one proves that $BM$ is globally
generated. These facts show that $\psi^*B$ is very ample on $S$:
the $B$ sections recover the embedded base $P$, while, on an open
where a section of $BM$ is nonzero, its corresponding section of
$\psi^*B$ is a unit multiple of the cyclic generator $z$.
They generate the local algebra of the degree-three cover and hence
give a closed immersion. Equivalently they separate length-two
subschemes whose image is a point as well as those separated on $P$.

Riemann--Roch and the same vanishing give the exact dimensions
\[
 h^0(P,B)=210,\quad h^0(P,BM)=186,\quad
 h^0(P,BM^2)=160,\quad h^0(S,Q^9)=556.
\]
For complete explicit defining coefficients, the ordered basis
$v_1,\ldots,v_{556}$ is the basis over $\mathbb F_2$ in
\texttt{code/A3\_Q9\_basis.json}, generated by
\texttt{code/certify\_A3.py}. Here is its finite mathematical
specification. For $i=0,1,2$ put $a=6-i$, $b=9+2i$, and order the
monomials $w^jx^ry^sz^i$ lexicographically by $(j,s,r)$, where
\[
 0\leq j\leq a,\quad s\in\{0,1\},\quad r\geq0,
 \qquad 2r+15s\leq2b+7a-j.
\]
Use the expansion $x=t^{-2}u$, $y=t^{-15}u^8$, where
$u=1+t^{15}u^{-7}$. Impose that, for each $0\leq h\leq a$,
every negative power of $t$ in
\[
 \sum_{j\geq h,r,s}c_{jrs}{j\choose h}
       t^{2b+7h-j-2r-15s}u^{r+8s}
\]
has coefficient zero. This is precisely the regularity condition
after $w_U=t^6w_V+t^{-1}$ and $z_U=t^4z_V$.
Order the equations by their first occurrence when the ordered
monomials are processed, and take the reduced free-variable kernel
basis with free variables in increasing order, as specified by the
included elimination routine. Concatenate the three bases in
increasing $i$. All sections are regular on the affine curve, on
the chart at $\infty$, and along the section $T$ by their indicated
projective powers. Independence and the Riemann--Roch dimensions prove
that they form the complete basis.

The exact calculation has respectively $357,291,230$ variables,
ranks $147,105,70$, and nullities $210,186,160$. Only four, four,
and three terms beyond the constant in the $t^{15}$ expansion can
contribute to the negative powers, so the calculation is finite and
exact. Its output and time are recorded in
\texttt{code/certify\_A3\_output.json}; the basis SHA256 is recorded
there too. These checks verify coefficient and transition equations,
not a sampled collection of points on the surface.

\paragraph{5. The cover making the first power very ample.}
Let
\[
 k=\overline{\mathbb F_2(\lambda_{ij}\mid1\leq i,j\leq556)},
 \qquad s_i=\sum_{j=1}^{556}\lambda_{ij}v_j.
\]
All varieties and bundles above are now extended to $k$. The matrix
$\lambda$ is invertible, so the $s_i$ are a basis of $H^0(S,Q^9)$.
Their zero divisors are smooth, pairwise transverse, and no three
meet. Here is the exact genericity argument. The smooth-divisor
Bertini theorem applies because $Q^9$ is very ample (Stacks,
Lemma~33.47.3, opened)~\cite{StacksBertini}. After selecting any number of such divisors,
apply the same theorem to their smooth components to make the next
divisor transverse, and avoid their finitely many pairwise
intersection points. Also avoid the linear span of the already
chosen sections. At each stage these are finitely many nonempty
open conditions. Their intersection is nonempty over an infinite
field. They define a nonempty open set of ordered bases; algebraic
independence of all $\lambda_{ij}$ places our basis in this set.

For each $s_i$ take the ninth-root algebra
$\bigoplus_{r=0}^{8}Q^{-r}$ with multiplication determined by
$s_i:Q^{-9}\to\OO_S$, and set
\begin{equation}\label{a:root-algebra}
 \mathcal A=\bigotimes_{i=1}^{556}
       \left(\bigoplus_{r=0}^{8}Q^{-r}\right),\qquad
 \rho:Y=\operatorname{Spec}_S\mathcal A\longrightarrow S.
\end{equation}
This describes the whole finite cover, not a choice of a connected
component. It has degree $9^{556}$. There are tautological sections
$t_i$ of $L=\rho^*Q$ with $t_i^9=\rho^*s_i$.

The surface $Y$ is smooth. At any point of $S$ at most two branch
divisors pass, with independent regular parameters as equations.
After an étale extension extracting ninth roots of units, the cover
is obtained by replacing each such parameter by its ninth power.
The other equations define étale extensions, since nine is
invertible. The resulting local rings are regular of dimension two.
Projectivity follows from finiteness. Moreover
\begin{equation}\label{a:root-directsum}
 \rho_*\OO_Y=
 \bigoplus_{0\leq r_1,\ldots,r_{556}\leq8}
                   Q^{-(r_1+\cdots+r_{556})}.
\end{equation}
Ampleness of $Q$ makes all negative-power global sections zero, so
$H^0(Y,\OO_Y)=k$. Thus $Y$ is connected, and its smoothness over the
algebraically closed field makes it integral.

We verify very ampleness of $L$ without conflating this cover with
the inverse image under a coordinate power map. Use the six pulled
back sections of $Q$ and the $556$ sections $t_i$. They generate $L$
because the old six already do. Let $u$ be one of the old sections.
The open $S_u=S\setminus V(u)$ is affine: it is the intersection of
the $Q^9$ embedding with the affine hyperplane chart $u^9\ne0$.
Since the $s_i$ form a complete basis, its coordinate algebra is
generated over $k$ by $s_i/u^9$. Over this open put $z_i=t_i/u$.
Then
\[
 \OO(Y_u)=\OO(S_u)[z_1,\ldots,z_{556}]
                   /(z_i^9-s_i/u^9)_{i=1}^{556}.
\]
This algebra is generated over $k$ by the $z_i$ alone, because their
ninth powers include all the displayed generators of $\OO(S_u)$.
Consequently the projective morphism defined by the old sections and
the $t_i$ is a closed immersion on each target chart $u\ne0$.
Those charts cover $Y$. The morphism is therefore an immersion,
and properness makes it a closed immersion globally. Hence $L$ is
very ample. Retaining the old sections in this argument also
separates simultaneous multiplication of all $t_i$ by a ninth root
of unity; no torsion ambiguity in the identity $L=\rho^*Q$ remains.

Finally tensor~\eqref{a:root-directsum} by $Q^{-1}$ and use
\eqref{a:Q-cohomology} and the higher negative-power vanishing.
Only the multi-index $(0,\ldots,0)$ contributes to $H^1$, giving
$h^1(Y,L^{-1})=1$. The same decomposition gives
$h^0(Y,L)=6+556=562$ and $h^1(Y,\OO_Y)=9+556=565$.
This completes all geometric and cohomological assertions in
Theorem~\ref{a:very-ample-example}.

\paragraph{Verification and attribution.}
Three other AI agents independently checked the root-cover argument,
the smoothness and connectedness assertions, the cohomology summands,
and the very-ampleness chart proof; their audit notes are retained.
The finite certificate is independent of those prose checks. The
construction uses the classical Raynaud mechanism, but the equations,
specified polarization and simultaneous-root proof above are supplied
in full here. Searches for a published theorem with this exact
very-ample negative-$H^1$ conclusion did not locate one during this
run. No priority claim is made.

\subsection{The stronger generic-kernel assertion of SA14-H1}
\begin{proposition}[R3u/SA14-H1: generic kernel is infinitesimal unipotent; Proved]\label{a:generic-unipotent}
Let $X$ be a normal connected projective variety of dimension at least
two over an algebraically closed field $k$ of characteristic $p>0$.
For a specified very ample sheaf $L$, let $Y$ be the geometric generic
hyperplane section. Then the kernel
\[
 G=\ker\bigl(\mathbf{Pic}^{\tau}_{X/k}\otimes_k\overline{k(|L|)}
          \longrightarrow\mathbf{Pic}^{\tau}_{Y/\overline{k(|L|)}}\bigr)
\]
is finite, connected, and unipotent.
\end{proposition}
\begin{proof}
The finiteness proof in Proposition~\ref{a:generic} applies to normal
$X$: $\mathbf{Pic}^{\tau}$ is proper of finite type, and a
positive-dimensional reduced connected kernel would contain
prime-to-$p$ torsion, contrary to the fundamental-group surjection.

The argument excluding geometric points also works for normal $X$,
with one precision concerning differentials. Associate to a
$p$-torsion line bundle its rational logarithmic differential in
$\Omega^1_{k(X)/k}$, viewed as a regular form on $X_{\mathrm{sm}}$.
If it is zero, the local units describing a trivialization of the
$p$th tensor power have $p$th roots in $k(X)$. Normality implies that
each root and its inverse are regular on the relevant trivializing
open, so the roots trivialize the line bundle. Thus this rational
$d\log$ is injective on $\Pic(X)[p]$.

The incidence argument is carried out over $X_{\mathrm{sm}}$ and the
open locus of hyperplanes smooth at the specified point. Its relative
cotangent sheaf is locally free. A form defined over $k$ that vanishes
on the geometric generic hyperplane therefore vanishes throughout
this integral open incidence locus. At a smooth point where the form
is nonzero, very ampleness lets the tangent hyperplane vary over all
codimension-one subspaces. Some such subspace is not annihilated by
the form, a contradiction. Triviality of the restricted line bundle
implies vanishing of its relative logarithmic differential on the
hyperplane: its global units are constants. As in
Proposition~\ref{a:generic}, every finite-order geometric point of the
Picard scheme is defined over $k$. The argument excludes $p$-torsion,
and taking powers excludes all $p$-power torsion. Consequently $G$
has just its identity as a geometric point and is connected.

It remains to prove unipotence, which does not follow from
connectedness alone. Raynaud's Proposition~6.2.1, p.~50, and
Proposition~6.2.3(1), p.~51, give a natural isomorphism
\[
 \Hom(\mu_p,\mathbf{Pic}_{Z/\Omega})
       \simeq H^1_{\text{ét}}(Z,\ZZ/p)
\]
for a proper scheme $Z$ over an algebraically closed field $\Omega$
\cite{Raynaud1970PicardDuality}. Apply it to $X_\Omega$ and $Y$.
The surjection $\pi_1(Y)\to\pi_1(X_\Omega)$ makes restriction
injective on these cohomology groups. Hence $\Hom(\mu_p,G)=0$.

Here is an elementary finite-group argument completing the proof.
Let $E$ be the maximal étale quotient of the Cartier dual $G^D$.
Its dual $E^D$ is a closed subgroup of the finite connected scheme
$G$, and hence is connected. Over $\Omega$, $E$ is a constant finite
abelian group; connectedness of its diagonalizable dual implies
that it is a $p$-group. If $E\ne0$, a quotient $E\to\ZZ/p$ dualizes
to a subgroup $\mu_p\subset E^D\subset G$, a contradiction.
Thus $G^D$ is connected, so its finite coordinate algebra is local
with residue field $\Omega$.

Finite-dimensional representations of $G$ are modules over
$\mathcal O(G)^*=\mathcal O(G^D)$. That local algebra has only the
trivial simple module. Every such representation therefore has a
finite filtration with trivial one-dimensional factors. Applying
this to the faithful regular representation embeds $G$ in a group
of upper unitriangular matrices. This is unipotence.
\end{proof}

For normal $X$, the tangent-kernel computation also remains valid.
The geometric generic hyperplane is geometrically integral, so
$H^0(X_\Omega,\mathcal O)=H^0(Y,\mathcal O)=\Omega$; the sequence
$0\to L^{-1}\to\mathcal O_X\to\mathcal O_Y\to0$ identifies the tangent
space of $G$ with $H^1(X,L^{-1})\otimes_k\Omega$. A finite connected
group scheme is nontrivial exactly when its tangent space is nonzero,
by Nakayama applied to its local coordinate algebra.

\subsection{A conditional source example has incompatible hypotheses}
\begin{proposition}[A-CZ346ii: incompatible source hypotheses; Proved]\label{a:cz-vacuity}
Let $i:X\to Y$ be a morphism of smooth schemes over a perfect field
$k$ of characteristic $p>0$. If
\[
 H^0(Y,\Omega^1_{Y/k})\longrightarrow H^0(X,\Omega^1_{X/k})
\]
is injective, then
$H^1_{\mathrm{fppf}}(Y,\alpha_p)\to
 H^1_{\mathrm{fppf}}(X,\alpha_p)$ is injective.
\end{proposition}
\begin{proof}
There is a natural injective additive map
\[
 H^1_{\mathrm{fppf}}(Z,\alpha_p)
       \longrightarrow H^0(Z,\Omega^1_{Z/k})
\]
for every smooth $k$-scheme $Z$; its scalar convention is Frobenius
semilinear, which does not affect injectivity. On affine opens,
the sequence $0\to\alpha_p\to\mathbf G_a\xrightarrow{F}\mathbf G_a\to0$
writes an $\alpha_p$-torsor as $z^p=a_j$. On overlaps the functions
$a_j$ differ by $p$th powers, so $da_j$ glue. If the glued form is
zero, smoothness over the perfect field implies locally
$a_j=b_j^p$. For instance, this follows first in the function field
from $\ker(d)=k(Z)^p$, and then in local rings from normality.
The torsor is therefore Zariski locally trivial. Its local
trivializations glue uniquely, since $\alpha_p$ has no nonzero
sections on a reduced scheme. This proves injectivity.
The construction commutes with pullback. Applying it to $X$ and
$Y$ proves the assertion.
\end{proof}

Cotner--Zavyalov, Example~3.3.6, supplies a nonzero kernel on
$H^1(\alpha_p)$, while the additional hypothesis of their
Example~3.4.6(ii), printed pp.~1867--1868, asks that $X$ also be
smooth and that restriction on global one-forms be an isomorphism
\cite{CotnerZavyalov2024}. The proposition shows that these
hypotheses cannot hold simultaneously. Thus that conditional
example gives no counterexample, or realizable remaining gate,
for injectivity on $\Pic[p]$. This observation concerns the
auxiliary example only, not their main Lefschetz theorems.

\subsection{A finite determinantal test for special hyperplanes}

\begin{proposition}[A2D; Proved]\label{a:determinantal}
Let $X/k$ be smooth projective and integral of dimension at least two,
where $k$ is algebraically closed of characteristic $p>0$.
Let $L$ be very ample, let $M\ne\OO_X$ satisfy $M^p\simeq\OO_X$,
and put
\[
 a=h^1(X,M\otimes L^{-1}),\qquad b=h^1(X,M),\qquad
 r=h^0(X,L)-1.
\]
On the open subset of $|L|=\mathbb P^r$ parametrizing smooth divisors,
the locus where $M|_Y$ is trivial is finite. It is the restriction of
the maximal-minor locus of the linear multiplication matrix
\[
 H^1(X,M\otimes L^{-1})\otimes\OO_{\mathbb P^r}(-1)
 \longrightarrow H^1(X,M)\otimes\OO_{\mathbb P^r}.
\]
We have $a\leq b$. If this finite locus is nonempty, then $a>0$ and
\begin{equation}\label{a:determinantal-bound}
 b-a\ \geq\ h^0(X,L)-2.
\end{equation}
In dimension at least three, it follows that restriction on geometric
Picard groups is injective for every smooth member of $|L|$ if, for
each nontrivial $M\in\Pic(X)[p]$, either $a=0$ or the reverse strict
inequality to \eqref{a:determinantal-bound} holds.
\end{proposition}

\begin{proof}
The injective logarithmic differential construction of
Proposition~\ref{a:generic} gives a nonzero form
$\eta=d\log(M)\in H^0(X,\Omega_X^1)$.
If $M|_Y$ is trivial, its restriction to $\Omega_Y^1$ is zero.
The conormal sequence and
$H^0(Y,\OO_Y(-Y))=0$ therefore show that
$\eta|_Y=0$ already in $\Omega_X^1|_Y$.
Every such $Y$ is consequently a divisorial component of the zero
scheme of $\eta$. A nonzero section of a vector bundle on an integral
noetherian scheme has only finitely many divisorial components in its
zero scheme. A smooth ample divisor here is connected and hence
integral. There are therefore only finitely many possible $Y$.

A nontrivial numerically trivial line bundle on a projective integral
scheme has no nonzero section: any nonzero effective divisor has
positive intersection with a suitable power of an ample divisor.
The same intersection argument gives $H^0(X,M\otimes L^{-1})=0$.
For a section $h$ cutting out a smooth $Y$, the divisor sequence gives
\[
 H^0(Y,M|_Y)=\ker\bigl(H^1(X,M\otimes L^{-1})
                  \xrightarrow{\,h\,}H^1(X,M)\bigr).
\]
Since $M|_Y$ is numerically trivial, its space of sections has dimension
one exactly when it is trivial, and dimension zero otherwise.
This proves the asserted determinantal description. There are smooth
members outside the finite exceptional locus, so a multiplication map
is injective and $a\leq b$.

At a bad smooth member, the matrix has rank exactly $a-1$.
If $a=0$ there is no such member. Otherwise choose an invertible
$(a-1)$-minor locally, allowing the empty minor when $a=1$.
Elementary row and column operations reduce the maximal-minor ideal
locally to $b-a+1$ functions, the entries of its remaining column.
The bad locus has dimension zero in the smooth-member open subset
of the regular variety $\mathbb P^r$. The height theorem for an ideal
generated by $b-a+1$ elements therefore gives $r\leq b-a+1$,
which is \eqref{a:determinantal-bound}.

In dimension at least three the point kernel is a finite $p$-group
by the already proved A2 reduction. Any nonzero such group contains
an element of order $p$. Applying the preceding exclusion to every
nonzero $p$-torsion line bundle proves the last assertion.
\end{proof}

This supplies a finite linear-algebra test for a fixed pair $(X,L)$
once its $p$-torsion classes and multiplication maps are computed.
The entries depend linearly on the coefficients of the hyperplane.
The full special-hyperplane problem is not settled: the inequality
can hold, and we have not proved that the resulting determinantal
points must lie in the discriminant of singular divisors.

\subsection{The leading Cartier term of a possible point kernel}
\paragraph{A2V (Proved necessary condition; not a solution of A2).}
Let $X/k$ be smooth, projective and integral over an algebraically
closed field of characteristic $p>0$, with $\dim X\geq2$.
Let $Y\subset X$ be a smooth ample divisor and put
$L=\mathcal O_Y(Y)$. Suppose a nonzero closed global one-form
$\eta$ vanishes along $Y$. Let $m\geq1$ be its exact vanishing
order along the integral divisor $Y$. Then $p\mid m$.
If in addition $C\eta=\eta$, write $m=pa$. There is a nonzero
element
\[
 H^0(Y,B_Y^1\otimes L^{-a})
  \simeq
 \ker\bigl(F:H^1(Y,L^{-a})\longrightarrow
                   H^1(Y,L^{-pa})\bigr),                 \tag{A2V.1}
\]
where $B_Y^1=\operatorname{im}(d:F_*\mathcal O_Y\to
F_*\Omega_Y^1)$ has its Frobenius-module structure.
The cohomology group in the Frobenius target is understood
with the corresponding Frobenius scalar structure; equivalently,
the map is $p$-semilinear on the usual $k$-vector spaces.
In particular a nonzero $p$-torsion line bundle on $X$ that
restricts trivially to $Y$ forces this obstruction for some $a\geq1$.

Here is the proof. The leading coefficient of $\eta$
is a nonzero section
\[
 \sigma\in H^0\bigl(Y,\Omega_X^1|_Y\otimes L^{-m}\bigr).
\]
Locally write $\eta=t^m\alpha$, where $t$ cuts out $Y$.
The equation $d\eta=0$ implies
$m\,dt\wedge\alpha=0$ modulo $t$. If $p\nmid m$, the
leading coefficient is purely conormal. The conormal sequence
would put it in
\[
 H^0(Y,L^{-m-1})=0,
\]
contrary to its definition. Thus $m=pa$.

Since $d(t^{pa})=0$, the local form $\alpha$ is closed.
Cartier's rule $C(f^p\alpha)=fC(\alpha)$ and $C\eta=\eta$
give
\[
 C\alpha=t^{(p-1)a}\alpha.
\]
Restriction to the smooth divisor commutes with Cartier, so
the tangential leading coefficient $\beta$ has $C_Y\beta=0$.
It is nonzero: otherwise $\sigma$ would again be a section of
$L^{-pa-1}$, which has no nonzero global section. When local
equations change by $t_i=u_{ij}t_j$, the coefficients change
by $\alpha_i=u_{ij}^{-pa}\alpha_j$. Because the multiplier is
a $p$th power, closedness and exactness are preserved, and the
Frobenius-module interpretation makes these precisely the
transition functions of $B_Y^1\otimes L^{-a}$. The exact
Cartier sequence identifies $\ker C_Y$ on closed forms with
$B_Y^1$.

Finally, tensoring the Frobenius sequence with $L^{-a}$ gives
\[
 0\longrightarrow L^{-a}\longrightarrow F_*L^{-pa}
  \longrightarrow B_Y^1\otimes L^{-a}\longrightarrow0.
\]
Both negative line bundles have no global section, so its
cohomology sequence proves (A2V.1). The two exact sequences
and projection-formula conventions can also be checked in
\cite[Lemmas 9.2 and 9.5, pp.~2106--2107]{DGJ_Langer2011}.
For a point kernel take $\eta=d\log(M)$. It is nonzero and
Cartier-fixed; triviality of $M|_Y$ first makes its tangential
restriction zero, and the conormal sequence with
$H^0(Y,L^{-1})=0$ makes its full restriction zero.

There is also a numerical exclusion that can be checked on the
ambient variety. Since $m\geq p$, this same nonzero form defines
an injection of sheaves
\[
 \mathcal O_X(pY)\hookrightarrow\Omega_X^1.
\]
It is injective because $X$ is integral and the map is nonzero
at its generic point. Taking global sections gives
\[
 h^0(X,\Omega_X^1)\ \geq\ h^0(X,\mathcal O_X(pY)).
                                                        \tag{A2V.2}
\]
Consequently, if $\dim X\geq3$, the reverse strict inequality
excludes every nontrivial point in the restriction kernel by
the already established finite-$p$-group reduction. This bound
does not assert that such an inequality holds for all $X$.

\paragraph{Why Cartier iteration alone stops here.}
For every prime $p$, the affine local model
\[
 U=\operatorname{Spec}k[t,x,(1+t^px)^{-1}],\qquad
 \eta=d\log(1+t^px)=\frac{t^p\,dx}{1+t^px}
\]
is closed and Cartier-fixed and has exact order $p$ along
$t=0$. Its tangential leading term is the nonzero exact form
$dx$. Explicitly,
\[
 \eta=\sum_{n\geq0}(-1)^n t^{p(n+1)}x^n\,dx.
\]
Cartier retains $n=pr+p-1$ and transforms those terms into
$(-1)^r t^{p(r+1)}x^r\,dx$, giving the same series.
Thus one cannot infer divisibility of the vanishing order by
$p^2$, or infinite-order vanishing, from Cartier fixedness
alone. This is an affine local model, not a projective ample
counterexample. Its exact finite coefficient identities for
$p=2,3,5,7$ are recorded in
\texttt{check\_cartier\_leading.py} and the saved JSON output.

\paragraph{A necessary singularity of the associated cover.}
Any actual point-kernel counterexample has a nonnormal
$\mu_p$-torsor along $Y$. Indeed, the torsor has local equation
$z^p=u$, with $u$ a unit and $du/u=\eta$. It is integral:
if $u$ were a $p$th power in the function field, the generic
torsor section would extend over the normal scheme $X$,
trivializing $M$. Since $\eta$ vanishes along $Y$, the
Jacobian criterion makes this hypersurface singular at the
generic point over $Y$. That is a codimension-one point.
The integral hypersurface is Cohen--Macaulay, but fails
regularity in codimension one, so it is not normal. In
particular constructions whose $\mu_p$-cover stays normal
cannot yield the required point kernel. This observation
does not rule out the nonnormal-cover route.

\paragraph{Exact remaining gate.}
The necessary Frobenius kernel on the smooth ample divisor
need not vanish in positive characteristic. Its existence
does not construct a Cartier-fixed form on an ambient smooth
projective variety, and no converse is proved here. The
remaining task is either to exclude that global extension or
to realize it with a fully specified smooth very ample pair.

\subsection{An exclusion from closed global one-forms}
\paragraph{A2W (Proved under the stated differential hypothesis).}
Let $X/k$ be smooth, projective and integral of dimension $n\geq2$,
where $k$ is algebraically closed of characteristic $p>0$.
Suppose every global regular one-form on $X$ is closed:
\[
 d:H^0(X,\Omega_X^1)\longrightarrow H^0(X,\Omega_X^2)
 \quad\hbox{is zero}.
\]
If $L$ is very ample, then
\[
 H^0(X,\Omega_X^1\otimes L^{-1})=0.                 \tag{A2W.1}
\]
Consequently, if $n\geq3$, restriction on geometric Picard
groups is injective for every smooth divisor in $|L|$, using
the established reduction of its kernel to a finite $p$-group.
This is a conditional theorem: the closedness hypothesis is
not asserted for all smooth projective varieties in positive
characteristic.

To prove (A2W.1), suppose that
$0\ne\zeta\in H^0(X,\Omega_X^1\otimes L^{-1})$.
Choose a nonzero $s_0\in H^0(X,L)$ and put $\eta=s_0\zeta$.
For every $s\in H^0(X,L)$, both $s\zeta$ and $\eta$ are
global one-forms and hence closed. On the open set $s_0\ne0$
we therefore have
\[
 0=d\bigl((s/s_0)\eta\bigr)=d(s/s_0)\wedge\eta.
\]
Because $L$ gives a closed embedding, its affine coordinate
ratios $s/s_0$ generate the function field $K=k(X)$ over $k$.
Their differentials consequently span the $n$-dimensional
$K$-vector space $\Omega_{K/k}^1$. For a nonzero one-form
$\eta_K$ on a vector space of dimension at least two, the
linear map $v\mapsto v\wedge\eta_K$ is not zero. The
displayed equality for a spanning set is a contradiction.

For the Picard consequence, a nonzero finite $p$-group kernel
would contain a line bundle $M$ of order $p$. The nonzero
logarithmic form $d\log(M)$ restricts to zero tangentially on
the smooth divisor $Y$. The conormal exact sequence and
$H^0(Y,\mathcal O_Y(-Y))=0$ then make its full restriction
zero, putting it in the group (A2W.1). This proves injectivity.

There is a quantitative necessary condition even without the
closedness hypothesis. If a nonzero closed global form $\eta$
vanishes along a very ample divisor $Y$, write $\eta=s_0\zeta$
as above. The differentials of the global forms $s\zeta$ are
$d(s/s_0)\wedge\eta$. Their values over $K$ span the image
of the wedge map by $\eta$, which has rank $n-1$. It follows
that
\[
 \dim_k\operatorname{im}\bigl(
  d:H^0(X,\Omega_X^1)\to H^0(X,\Omega_X^2)\bigr)
  \geq n-1.                                         \tag{A2W.2}
\]
Thus any hypothetical point-kernel example in dimension
$n\geq3$ must have at least $n-1$ independent nonzero
differentials of global one-forms. In particular it must
have a nonzero Hodge-to-de Rham differential $d_1^{1,0}$.
The argument uses only closedness of the original logarithmic
form for this last necessary condition.

\subsection{Two target-side exclusions for a point kernel}
\paragraph{A2Y (Proved under either stated target hypothesis).}
Let $X/k$ be smooth, projective and integral of dimension
$n\geq3$ over an algebraically closed field of characteristic
$p>0$, and let $Y\subset X$ be a smooth very ample divisor.
Restriction
\[
 \operatorname{Pic}(X)\longrightarrow\operatorname{Pic}(Y)
\]
is injective on geometric line-bundle classes if either
\begin{enumerate}
\item Frobenius is injective on $H^1(Y,\mathcal O_Y)$; or
\item the smooth variety $Y$ has a smooth lifting to $W_2(k)$.
\end{enumerate}
No lifting of $X$ is assumed in the second assertion, and no
lifting of the very ample line bundle on $Y$ is required.
The first assertion includes $H^1(Y,\mathcal O_Y)=0$.
Frobenius injectivity here is the precise cohomological
hypothesis; it is not replaced by a claim about the reduced
Picard variety alone.

Put $L=\mathcal O_Y(Y)$, which is very ample on $Y$.
If the point kernel were nonzero, the finite-$p$-group
reduction gives a nontrivial class of order $p$. By A2V,
there would then be an integer $a\geq1$ and a nonzero class
\[
 \alpha\in H^1(Y,L^{-a}),\qquad F\alpha=0
       \quad\hbox{in }H^1(Y,L^{-pa}).                  \tag{A2Y.1}
\]
We exclude this class under each hypothesis.

First suppose Frobenius is injective on $H^1(Y,\mathcal O_Y)$.
Since $\dim Y\geq2$ and $L^a$ is very ample, Bertini gives
a smooth connected divisor $E\in|L^a|$, cut out by a section
$s$. Projectivity and geometric connectedness give
$H^0(Y,\mathcal O_Y)=H^0(E,\mathcal O_E)=k$. The exact
sequence
\[
 0\longrightarrow L^{-a}\xrightarrow{\,s\,}\mathcal O_Y
  \longrightarrow\mathcal O_E\longrightarrow0
\]
therefore makes multiplication by $s$ an injection
\[
 j_s:H^1(Y,L^{-a})\hookrightarrow H^1(Y,\mathcal O_Y).
\]
The Frobenius square commutes:
\[
\begin{array}{ccc}
 H^1(Y,L^{-a})&\xrightarrow{\ F\ }&H^1(Y,L^{-pa})\\
 {\scriptstyle j_s}\downarrow&&\downarrow{\scriptstyle j_{s^p}}\\
 H^1(Y,\mathcal O_Y)&\xrightarrow{\ F\ }&H^1(Y,\mathcal O_Y).
\end{array}
\]
Its horizontal arrows are $p$-semilinear, or linear when
their targets are given the Frobenius scalar structure.
Equation (A2Y.1) implies $F(j_s\alpha)=0$. The assumed
injectivity of the bottom arrow and then the injectivity of
$j_s$ give $\alpha=0$, a contradiction. No assertion that
$j_{s^p}$ is injective is needed.

Now suppose that $Y$ lifts smoothly to $W_2(k)$. Apply
\cite[Section 11 hypothesis and Corollary 11.2, p.~2116]{DGJ_Langer2011}
with variety $Y$, numerically flat bundle $\mathcal O_Y$,
ample divisor corresponding to $L^a$, and indices $i=0,j=1$.
Its required inequality is
\[
 1<\min(p,\dim Y),
\]
which holds for every prime $p$ and every $\dim Y\geq2$.
It yields $H^1(Y,L^{-a})=0$, again contradicting (A2Y.1).
The proof of that corollary uses the truncated
Deligne--Illusie inequality only in total degree one and
iterates Frobenius until Serre vanishing applies. It does
not require $\dim Y<p$.

\emph{Remaining gate.}
A counterexample to point injectivity must therefore have
noninjective Frobenius on $H^1(Y,\mathcal O_Y)$ and a target
$Y$ that admits no smooth lifting to $W_2(k)$, as well as
the independent ambient restrictions of A2V and A2W. These
are necessary conditions, not an existence construction.

\subsection{Targets covered by recent vanishing theorems}
\begin{proposition}[A1Q: quasi-F-split and Fano targets; Known deduction]
\label{a:known-targets}
Let $X$ be a smooth connected projective $n$-fold over an algebraically
closed field of characteristic $p>0$, and let $Y$ be a smooth effective
ample Cartier divisor. The following statements hold for this specified
divisor, and hence for a smooth hyperplane of a specified embedding.
\begin{enumerate}
\item If $Y$ is quasi-F-split and $n\geq3$, then
$\Pic(X)\to\Pic(Y)$ is injective. If also $n\geq4$, it is an isomorphism.
\item If $n=4$ and $Y$ is Fano, then
$\Pic(X)\to\Pic(Y)$ is an isomorphism, including in characteristic two.
\end{enumerate}
Here quasi-F-split means that, for some integer $r\geq1$, the natural
$\OO_{Y^{(1)}}$-linear map
$\OO_{Y^{(1)}}\to F_*W_r(\OO_Y)/p$ has an
$\OO_{Y^{(1)}}$-linear retraction, with the Frobenius module structure
as in \cite[\S1, p.~1]{RootPetrov2025}.
\end{proposition}
\begin{proof}
Put $d=\dim Y=n-1$ and $A=\OO_Y(Y)$, which is ample.
Petrov's Corollary~4.10(2), printed p.~16 of the opened version,
gives
\[
 H^{d-i}(Y,\omega_Y\otimes A^m)=0
       \qquad (m>0,\ 0<i<d)
\]
in the first case: use differential-form degree $d$, cohomological
degree $d-i$, and ample bundle $A^m$. The sum of the two degrees is
$2d-i>d$. Serre duality identifies the dual of the displayed group
with $H^i(Y,A^{-m})$. Thus all the negative first cohomology groups
vanish for $n\geq3$, and all the negative second groups vanish for
$n\geq4$. Apply Corollary~\ref{a:vanishing}.

For the second case, Kawakami--Tanaka prove Akizuki--Nakano vanishing
for every smooth Fano threefold in every positive characteristic:
\cite[Theorem~7.2, p.~56]{RootKawakamiTanakaII2025}, together with
\cite[Theorem~6.3, p.~50]{RootKawakamiTanakaI2026} for the case
$\Pic(Y)=\mathbf ZK_Y$ with very ample anticanonical bundle.
Take differential-form degree zero, cohomological degree $i=1,2$,
and ample divisor corresponding to $A^m$. Both sums of degrees are
strictly less than three, so $H^i(Y,A^{-m})=0$ for every $m>0$.
Corollary~\ref{a:vanishing} again applies. Neither proof requires
that the ambient variety lift, that the given divisor lift, or that
$p$ exceed its dimension.
\end{proof}

These are target hypotheses. They are not properties established for
an arbitrary smooth ample divisor. In particular this proposition
does not settle A1 or A2. The cited papers also distinguish Fano
threefolds from quasi-F-split varieties: some smooth Fano threefolds
are not quasi-F-split, so the second assertion is not redundant.

\subsection{An adjunction obstruction to a separable-cover construction}
\paragraph{A2Z (Proved exclusion of the stated construction).}
Let $k$ be algebraically closed of characteristic $p>0$, and let
$\pi:X\to B$ be a finite, flat, generically separable morphism of
degree $p$ between smooth projective integral varieties of dimension
$n\geq2$. Let $C\subset B$ be a smooth ample integral divisor, and
assume that its scheme-theoretic inverse image
\[
 Y=\pi^*C
\]
is smooth and integral. Then
\[
 \operatorname{Pic}(C)[p](k)\longrightarrow
 \operatorname{Pic}(Y)[p](k)
\]
is injective. In particular this construction cannot make a
nontrivial $M|_C$, for $M\in\operatorname{Pic}(B)[p](k)$, disappear
on $Y$. This is a restriction on the proposed construction; it does
not prove the general point-injectivity assertion A2.

Suppose that $N$ is a nontrivial order-$p$ line bundle on $C$ and
$\rho^*N$ is trivial, where $\rho=\pi|_Y$. Choose $N^p\simeq
\mathcal O_C$, and let $\tau:T\to C$ be the associated
$\mu_p$-torsor. It is integral. Indeed, if its generic equation
$z^p=u$ split in $k(C)$, the rational root would give a rational
trivialization of $N$ with zero divisor, hence a global
trivialization because $C$ is normal. Its degree is $p$.
A trivialization of $\rho^*N$ can be made compatible with the
chosen $p$th-power trivialization by rescaling by a constant:
$H^0(Y,\mathcal O_Y^*)=k^*$ and $k$ is algebraically closed.
It therefore gives a factorization
\[
 Y\xrightarrow{\nu}T\xrightarrow{\tau}C.
\]
Since both $\rho$ and $\tau$ have degree $p$, $\nu$ is finite
birational. Thus $Y$ is the normalization of $T$. In particular
$\rho$ is purely inseparable.

The torsor $T$ is Gorenstein and
\[
 \omega_T\simeq\tau^*\omega_C.                         \tag{A2Z.1}
\]
For an explicit check, realize it inside the total space of the
line bundle $N$ by its equation $z^p=u$. Adjunction gives
$\omega_T\simeq\tau^*(\omega_C\otimes N^{p-1})$; the
nowhere-zero tautological root trivializes $\tau^*N$, giving
(A2Z.1). This argument applies when $T$ is nonnormal.

Write $\mathfrak c$ for the conductor, viewed as an ideal of
$\mathcal O_Y$. Finite duality identifies
\[
 \omega_Y\simeq\mathfrak c\otimes\nu^*\omega_T.
\]
Here the conductor identification is literal: on an affine
birational extension $A\subset A'$ inside a common fraction field,
\[
 \operatorname{Hom}_A(A',A)
   =\{q:qA'\subset A\}
   =\operatorname{Ann}_A(A'/A).
\]
Both canonical sheaves are invertible, so $\mathfrak c$ is an
invertible nonzero ideal on $Y$. Consequently it is
$\mathcal O_Y(-D)$ for an effective Cartier divisor $D$, and
\[
 \omega_Y\simeq\rho^*\omega_C\otimes\mathcal O_Y(-D).
                                                               \tag{A2Z.2}
\]
The finite-duality input is the $R\operatorname{Hom}$ description
of the dualizing complex under a finite map
\cite[Lemma 47.15.8, with the finite-module description in Section 49.5]{DGJ_StacksFiniteDuality};
taking the common lowest cohomological degree gives the displayed
module identity. The divisor adjunction used here also follows
from \cite[Lemma 48.14.1, Tag 0B4B]{DGJ_StacksAdjunction}.

On the other hand, the determinant of the generically invertible
map $d\pi:\pi^*\Omega_B^1\to\Omega_X^1$ is a nonzero regular
section of $\omega_X\otimes(\pi^*\omega_B)^{-1}$. Let $R$ be
its effective divisor. Hence
\[
 \omega_X\simeq\pi^*\omega_B\otimes\mathcal O_X(R).
                                                               \tag{A2Z.3}
\]
The divisor $R$ contains $Y$. At its generic point the normal
map in the conormal sequences for $Y\subset X$ and $C\subset B$
is an isomorphism, because $\pi^*C=Y$ scheme-theoretically,
whereas the tangential determinant vanishes since $\rho$ is
purely inseparable of degree $p$. Thus the full determinant
vanishes at that generic point.
Applying smooth-divisor adjunction to (A2Z.3) gives
\[
 \omega_Y\simeq\rho^*\omega_C\otimes
                  \mathcal O_X(R)|_Y.                         \tag{A2Z.4}
\]
Comparison with (A2Z.2) says
$\mathcal O_X(R)|_Y\otimes\mathcal O_Y(D)\simeq\mathcal O_Y$.
Intersecting with $\mathcal O_X(Y)|_Y^{\,n-2}$ gives
\[
 0=R\cdot Y^{n-1}+D\cdot (Y|_Y)^{n-2}>0,
\]
a contradiction: $Y$ is ample because it is a finite pullback
of $C$, $R$ is a nonzero effective divisor (indeed $R\geq Y$),
and $D$ is effective. The first summand is strictly positive and
the second is nonnegative. This completes the proof.

\emph{Scope of the obstruction.}
A local separable equation such as
$z^p-s^{p-1}z=b$ can have a smooth total space and purely
inseparable residue extension at $s=0$. The proposition explains
why these local equations cannot be glued into the stated
projective construction with a smooth ample full inverse image
that kills a nontrivial $p$-torsion class. Allowing the
$\mu_p$-torsor on $C$ to be nonnormal does not evade the argument:
its conductor contributes the additional nonnegative term above.
The degree-$p$ hypothesis and the equality $\pi^*C=Y$ are retained;
no assertion about arbitrary higher-degree covers or a single
component of a reducible inverse image is made.

\subsection{An inseparable multisection construction is also excluded}
\begin{proposition}[A2P: the relative Frobenius multisection obstruction]
Let $S$ be a smooth projective integral variety of positive dimension
over an algebraically closed field of characteristic $p>0$.
Use the quotient convention for $P=\mathbf P_S(E)$, where $E$ has
rank two. Let $F_S^*E\twoheadrightarrow N$ be a line quotient.
The relative Frobenius
$P\to\mathbf P_S(F_S^*E)$ has degree $p$ on each projective-line
fibre. Let $Y\subset P$ be the inverse image of the section
defined by that quotient, and suppose $Y$ is smooth and integral.
If $Y$ is ample, then
\[
       \Pic(S)[p](k)\longrightarrow\Pic(Y)[p](k)
\]
is injective. Thus this construction cannot give an A2 example
by normalizing a nontrivial $\mu_p$-torsor on $S$.
\end{proposition}
\begin{proof}
Write $\rho:Y\to S$ and $K=\OO_P(1)|_Y$. The morphism $\rho$
is finite flat and purely inseparable of degree $p$. The section
of $\mathbf P_S(F_S^*E)$ is cut out by its kernel line bundle
$(\det E)^p\otimes N^{-1}$. Pulling back its divisor gives
\[
 \OO_P(Y)=\OO_P(p)\otimes
                 \pi^*((\det E)^{-p}\otimes N),
 \qquad K^p=\rho^*N.                                  \tag{A2P.1}
\]
Put $Q=K^2\otimes\rho^*(\det E)^{-1}$. The projective-bundle
canonical formula and smooth-divisor adjunction give
\[
 \OO_P(Y)|_Y=Q^p,\qquad
 \omega_Y\otimes(\rho^*\omega_S)^{-1}=Q^{p-1}.          \tag{A2P.2}
\]
Indeed $\omega_P=\OO_P(-2)\otimes\pi^*(\omega_S\otimes\det E)$;
substituting (A2P.1) gives exponents $2p-2$ on $K$ and $1-p$
on $\rho^*\det E$, exactly as in (A2P.2).

Suppose a nontrivial $M\in\Pic(S)[p](k)$ becomes trivial on $Y$.
As in A2Z, its integral $\mu_p$-torsor $T\to S$ has degree $p$.
A compatible trivialization of $\rho^*M$ gives a finite birational
map $Y\to T$, so $Y$ is its normalization. The Gorenstein torsor
has $\omega_T\simeq\omega_S|_T$, and finite duality gives
\[
        \omega_Y\otimes(\rho^*\omega_S)^{-1}
                    \simeq\OO_Y(-D)                         \tag{A2P.3}
\]
for an effective conductor divisor $D$. The precise conductor
argument, including nonnormal $T$, was proved in A2Z using
\cite[Lemma 47.15.8]{DGJ_StacksFiniteDuality}.
Equations (A2P.2)--(A2P.3) imply $Q^{p-1}\simeq\OO_Y(-D)$.
If $Y$ is ample in $P$, its normal line $Q^p$ is ample, hence
$Q$ is ample. An ample positive power cannot be the negative of
an effective divisor on the positive-dimensional projective
variety $Y$: intersection with any ample class of complementary
degree is positive on the former and nonpositive on the latter.
This contradiction proves injectivity.
\end{proof}

This is an exclusion for the specified relative Frobenius
multisection. It neither assumes a disjoint infinity section nor
classifies arbitrary smooth ample divisors in projective bundles.
The unrestricted A2 question remains undecided.

